\section{Introduction}
\label{sec:introduction}

Ensemble Kalman filters require the inverse of the background error covariance matrix, but in most operational settings the ensemble is far too small to estimate that matrix reliably. When the state dimension $n$ exceeds the ensemble size $N_e$, the sample covariance is rank-deficient and cannot be inverted at all~\cite{houtekamer1998data}; even when $N_e > n$, sampling noise introduces spurious correlations that contaminate the precision matrix~\cite{anderson2001ensemble}. The same problem appears in cosmological data analysis, where the number of mock catalogs available to estimate power spectrum covariances is often comparable to the data dimension~\cite{hamilton2006measuring}.

Shrinkage estimation addresses this by combining the noisy sample estimate with a structured, low-variance target matrix, accepting some bias in exchange for reduced mean squared error~\cite{ledoit2004well}. Most shrinkage work operates in covariance space: Ledoit and Wolf~\cite{ledoit2004well} derived an optimal intensity for shrinking toward a scaled identity, Pope and Szapudi~\cite{pope2008shrinkage} brought this idea to cosmology, and Chen et al.~\cite{chen2010shrinkage} proposed oracle-approximating algorithms that improve on Ledoit-Wolf in high dimensions. Bodnar et al.~\cite{bodnar2016direct} took a different route, shrinking the precision matrix directly rather than first shrinking the covariance and then inverting. Looijmans et al.~\cite{looijmans2024comparison} applied this direct precision framework to cosmological data, comparing several target choices and showing that domain-informed targets outperform generic ones when the true precision is near-diagonal.

Two gaps remain. First, Looijmans et al.\ used these estimators for standalone precision matrix estimation, not inside a sequential filter. It is unclear whether an estimator that performs well in a one-shot setting will also perform well when applied repeatedly across hundreds of assimilation cycles, where errors in the precision estimate feed back through the forecast model. Second, the direct precision targets in the literature are either generic (identity, eigenvalues) or domain-specific (cosmological), with no middle ground for practitioners who lack prior knowledge of the precision structure but want something better than the identity.

We address both gaps. We embed the Bodnar et al.\ direct precision shrinkage framework inside the dual formulation of the Ensemble Kalman Filter, recomputing the shrinkage coefficients at each assimilation cycle. We propose a scaled identity target that uses empirical per-variable variances as a domain-agnostic alternative to the cosmological target. And we test all methods---three targets from Looijmans et al., our scaled variant, NERCOME~\cite{joachimi2017non}, and Ledoit-Wolf---on two systems with opposite precision structure: Lorenz~96~\cite{lorenz1996predictability} ($n = 40$, dense off-diagonal precision) and BOSS DR12 cosmological mock data~\cite{looijmans2024comparison} ($d = 18$, near-diagonal precision). The experiments show that direct precision shrinkage performance depends on how well the target matches the true precision structure, and that the estimator producing the worst standalone precision estimates (Ledoit-Wolf) yields the best EnKF filtering error, because its low variance pays off over repeated assimilation cycles. All implementations are publicly available in the TEDA framework~\cite{nino2022teda}.
